% !TEX root = ../dissertationOG.tex
% NOTE (2026-10-23): this Korean version lags the English revision: FlexiSim positioning and the four evidence questions (image information -> control paths -> physical retention -> generalization) are not yet translated here.
% 한국어판 4장 (Chapter4.tex의 번역, 2026-10-22 개정 3판: 세부는 부록으로). 한글 조판에 kotex(pdfLaTeX/XeLaTeX/LuaLaTeX)
% 또는 xeCJK(XeLaTeX)가 필요하다. 레이블은 영어판과 함께 쓸 수 있도록 모두 _ko 접미사를 붙였다. 그림 파일(figs/ch4/...)은
% 영어판과 공유한다. 3장 참조(\ref{CH3_SHAPEMORPHING})는 영어판 3장의 레이블을 그대로 가리킨다. 이 장의 세부(상수,
% 단계별 갱신, 항의 식, 선택 규칙, 여러 몸체 접촉, 구현, 척도, 긴 표)는 한국어 부록 Chapter4_appendix_ko.tex에 있다
% (영어판: Chapter4_appendix.tex, 학위논문 빌드: Sections/Appendix.tex). 인용 키는 영어판 docs/Chapter4.tex와 같다.

\clearpage % clear the prior chapter's page
\chapter{미분 가능한 물리 기반의 렌더 유도 체적 모핑}\label{CH4_PhysMorphGS_ko}
% 이 장이 쓰는 패키지: amsmath, graphicx, booktabs, array, tikz, algorithm + algpseudocode, natbib(\citep), 그리고 kotex 또는 xeCJK.
% 아래 TikZ 라이브러리는 표준 라이브러리다. 장 파일에서 그룹 밖에 두면 안전하지만, 학위논문 preamble에 두는 편이 낫다.
\usetikzlibrary{positioning,arrows.meta,fit,calc}

% ============================================================================
\section{개요}
\label{ch4:sec_overview_ko}
% ============================================================================

\ref{CH3_SHAPEMORPHING}장에서는 미분 가능한 MLS-MPM 시뮬레이션에서 입자별 제어 변형 구배를 최적화하여 한 형상을
다른 형상으로 모핑하였다. 제어는 내부 응력으로만 물질에 들어가므로 모핑은 질량과 운동량을 보존하고 위상을 바꿀 수 있다.
그러나 그 목적 함수는 배경 격자 위의 질량 손실이었고, 겉모습은 입자를 메시로 만들고 다듬어 렌더링할 때 나중에
만들어졌다. 관찰자가 보는 것, 즉 몸체의 실루엣도 표면의 셰이딩도 제어에 참여하지 않았다.

이 장에서는 그림을 제어의 일부로 만든다. 같은 종류의 물리적 몸체가 목표와 \emph{같아 보이도록} 최적화한다. 몸체 자신의
입자로 만든 표면에서 계산한 미분 가능한 렌더링 항이 제어를 통해 시뮬레이션을 이끌고, 제어는 여전히 응력과 표면의 강체
성분 없는 오프셋으로만 물질에 들어간다. 세 가지 관찰이 방법을 결정한다. 아직 밀고 있는 동안 평가되는 제어는 손을 놓으면
되튀어 돌아가는 상태를 얻을 수 있다. 그림은 MPM 셀보다 작은 픽셀 척도의 변화를 요구하는데, 격자를 통해서만 작용하는
제어는 그 요구를 한 셀의 많은 입자에 걸쳐 평균해 버린다. 그리고 렌더링한 입자 구름의 픽셀은 관찰자가 보는 표면의 픽셀이
아니다. 우리 파이프라인은 각각에 다음과 같이 답한다.
\begin{itemize}
  \item \textbf{안정화 창(settled window).} 몸체를 창 단위로 최적화하고, 각 창은 제어를 풀고 몸체가 가라앉은 뒤에만
        평가한다(\ref{ch4:sec_method_ko}절).
  \item \textbf{두 척도의 두 제어.} 상한이 있는 응력 증분 $\delta\mathbf{F}_c$는 \ref{CH3_SHAPEMORPHING}장처럼 격자를 통해
        덩어리를 옮기고, 바깥 입자층의 강체 성분 없는 법선 오프셋 $u$는 격자 아래에서 표면을 다듬는다
        (\ref{ch4:sec_controls_ko}절).
  \item \textbf{부피를 보존하는 몸체.} 입자는 자신의 운동이 만든 부피를 운반한다(\ref{ch4:sec_body_ko}절).
  \item \textbf{몸체 외피(exterior) 위의 렌더 항.} 입자 자신의 장이 만드는 표면에서 렌더 항을 읽고, 수송 기반 물리 목적
        함수와 함께 가중한다(\ref{ch4:sec_objective_ko}절).
\end{itemize}

이 장 전체를 관통하는 발견이 하나 있으며, 처음부터 밝혀 둔다. 렌더 항이 하는 일은 그 가중치를 정하는 방식에 달려 있다.
기본 규칙에서는 가중치가 물리 기울기를 따른다. 그러면 렌더는 몸체가 도착하는 동안 물질의 방향을 바꾸고, 도착한 뒤에는 얇은
부분의 테두리를 다듬으며, 끝은 우리가 읽은 모든 표준 척도에서 정확한 목표에 더 가깝지만 그 차이는 입자 간격의 일부이다.
도착 뒤에도 가중치를 유지하면 렌더는 그림만 담고 있는 것을 쓴다. 3D 목표가 매끈한 판인 몸체에 목표 그림에만 있는 글자를
새기고, 물리 목적 함수가 볼 수 없는 배치대로 색 있는 물질을 그림이 칠한 곳으로 옮긴다(\ref{ch4:sec_write_ko}절). 그 권한에는
대가가 있다. 실행이 멈출 때 몸체가 아직 움직이며, 권한과 정지를 함께 지키는 가중치는 \ref{ch4:sec_costs_ko}절의 열린 질문이다.
각 부분을 렌더 가중치만 0으로 둔 같은 실행인 $\lambda=0$ 쌍둥이(twin)와 비교하여 잰다. 상수, 단계별 갱신, 긴 표는 부록~\ref{app:ch4_ko}에 모았다.

\begin{figure*}[t]
    \centering
    \resizebox{\textwidth}{!}{\input{figs/ch4/pipeline_overview.tikz}}
    \caption{파이프라인 개요. \textmd{각 창은 구동 단계와 해제 단계를 시뮬레이션하고, 해제된 끝을 물리 항과 18개 시점에서
    본 몸체 외피의 실루엣·셰이딩으로 평가하며, 하나의 수반으로 두 제어를 갱신한다(부록~\ref{ch4:sec_cleanup_ko}의 두 정리
    항은 생략; 도식의 문구는 영어판과 공유).}}
    \label{fig:overview_ko}
\end{figure*}

% ============================================================================
\section{시뮬레이션되는 몸체}
\label{ch4:sec_body_ko}
% ============================================================================

\ref{CH3_SHAPEMORPHING}장의 MLS-MPM 시뮬레이터를 3차 B-스플라인 가중치, APIC 전달~\citep{jiang2016material}, 고정 공회전
탄성 응답과 함께 쓴다. 원천(구)과 목표 메시는 각각 같은 질량의 입자 $N$개로 샘플링하고, 원천은 목표의 부피에 맞추며,
입자당 부피의 모서리 길이 $a=(V/N)^{1/3}$을 \emph{샘플링 피치}라 부른다. MPM 셀은 $N$이 아니라 형상으로 정하므로, 입자
300k에서 한 셀에는 약 183개의 입자가 들어간다. 이것이 그림이 요구하는 것의 많은 부분이 격자 아래에 놓이는 이유이다. 구동
단계에서 제어는 탄성 변형에 들어가고,
\begin{equation}
\label{eq:control_F_ko}
\mathbf{F}^{\mathrm{E}}_p=\big(\mathbf{F}_p+\delta\mathbf{F}_{c,p}\big)\big(\mathbf{F}^{\mathrm{P}}_p\big)^{-1},
\qquad
\boldsymbol{\tau}_p=\mathbf{P}\big(\mathbf{F}^{\mathrm{E}}_p\big)\,\mathbf{F}^{\mathrm{E}\top}_p ,
\end{equation}
Kirchhoff 응력 $\boldsymbol{\tau}_p$는 \ref{CH3_SHAPEMORPHING}장처럼 입자-격자 전달에 들어간다.

세 가지 추가가 몸체의 부피와 표면을 정돈한다(부록~\ref{app:ch4_body_ko}). 각 입자는 자신의 운동이 만든 부피를 운반하므로,
\ref{CH3_SHAPEMORPHING}장에서 쓴 변형 구배의 평활이 더 이상 이동 중인 얇은 부분을 팽창시키지 않는다. 매 확정 때 소성 동화가
부피를 포함한 탄성 늘어남의 절반을 정지 형상으로 바꾸므로, 모핑의 끝은 몸체 자신의 정지 상태이다. 그리고 세 가지 위치
갱신이 격자 아래에서 작용한다. 너무 가까운 입자를 크기가 같고 방향이 반대인 밀어냄으로 떼어 놓는 최소 간격, 바깥층을 목표의
부조 쪽으로 이완하는 강체 성분 없는 이완, 떨어진 조각을 이웃과 함께 붙잡아 두는 결합이다.

% ============================================================================
\section{응력 제어와 표면 제어}
\label{ch4:sec_controls_ko}
% ============================================================================

\ref{CH3_SHAPEMORPHING}장처럼 정해진 궤적이나 외력이 아니라 물질의 내부 변형으로 모핑을 이끈다. 그러나 렌더링 항은 두
척도의 변화를 요구한다. 몸체의 덩어리는 목표까지 이동해야 하고 격자가 이를 운반할 수 있지만, 표면은 목표의 윤곽과 부조를
가져야 하며 그중 상당 부분은 MPM 셀 하나보다 작다. 그래서 최적화기에 두 제어를 준다.

\subsection{응력 증분}
$T$개의 구동 단계마다 모든 입자는 $3\times3$ 증분 $\delta\mathbf{F}^{n}_{c,p}$, 곧 \ref{CH3_SHAPEMORPHING}장의 제어 변형
구배를 가지며, 입자와 단계마다 상한을 둔다. 이 제어는 식~\eqref{eq:control_F_ko}를 통해서만 들어간다. 응력은
$\boldsymbol{\tau}_p(\mathbf{x}_i-\mathbf{x}_p)$로 격자에 들어가고 B-스플라인 가중치는
$\sum_i w_{ip}(\mathbf{x}_i-\mathbf{x}_p)=\mathbf{0}$을 만족하므로, 제어는 전체 선운동량을 바꿀 수 없다. 또한
$\boldsymbol{\tau}_p$가 대칭이므로 APIC는 각운동량도 보존한다~\citep{jiang2016material}. 최적화된 외력으로 몸체를 움직이는
방법도 생각할 수 있다. \ref{ch4:sec_force_ko}절에서 그렇게 바꾸고 무엇이 바뀌는지 잰다.

\subsection{바깥층 오프셋}
바깥층의 모든 입자는 창마다 스칼라 하나 $u_p$를 가진다. 이는 법선 방향의 변위로, 구동 단계에 걸쳐 $T$등분하여 적용한다.
\begin{equation}
\label{ch4:eq_u_ko}
\mathbf{x}_p\leftarrow\mathbf{x}_p+\Big(\tfrac{1}{T}\,g_p\,u_p+s'_p\Big)\mathbf{n}_p ,
\end{equation}
여기서 $s'_p$는 부록~\ref{ch4:sec_below_grid_ko}의 이완 단계, $g_p\in\{0,1\}$은 관문이다. 수반은 물리 응답에 항등 연산을
곱한 것이므로, 렌더 기울기는 격자의 저역 통과 없이 $u$에 닿는다. 세 규칙이 $u$를 손 편집이 아닌 표면 제어로 유지한다. $u$는
남은 수송이 MPM 셀 하나 이내인 입자에만 작용하므로, 물질을 나르는 것이 아니라 도착하는 표면을 다듬는다. 최적화기의 매 단계
뒤 바깥층의 여섯 강체 모드 성분을 제거하므로, 몸체를 어느 축을 따라서도, 어느 축에 대해서도 움직이지 않는다. 그리고 한
창에 한 간격으로 크기를 제한한다. $u$가 옮긴 것을 물리가 붙잡는지는 $u$가 작용하지 않는 창의 해제 절반이 결정한다.

% ============================================================================
\section{목적 함수}
\label{ch4:sec_objective_ko}
% ============================================================================

목적 함수는 세 묶음, 여덟 개의 항으로 이루어지며,
\begin{equation}
\label{ch4:eq_objective_ko}
\mathcal{L}=\underbrace{\sigma\big(\mathcal{L}_{\mathrm{OT}}+\mathcal{L}_{\mathrm{prox}}+\mathcal{L}_{\mathrm{drift}}\big)
 +\mathcal{L}_{\mathrm{rel}}}_{\mathcal{L}_{\mathrm{phys}}}
 \;+\;w_c\underbrace{\big(\mathcal{L}_{\mathrm{spray}}+\mathcal{L}_{\mathrm{near}}\big)}_{\mathcal{L}_{\mathrm{clean}}}
 \;+\;\lambda\underbrace{\big(\mathcal{L}_{\mathrm{sil}}+\mathcal{L}_{\mathrm{shade}}\big)}_{\mathcal{L}_{\mathrm{render}}} ,
\end{equation}
각 항의 식은 부록~\ref{app:ch4_terms_ko}에 있다. 19개 메시 갤러리와 입자 300k에서 항을 끄면 측정량이 바뀌는 경우에만 그
항을 남겼다. 기하 항과 안정 항에는 따로 가중치가 없다. $\sigma$는 원천에서 한 번 잰 상수 하나로, 기하 에너지에
\ref{CH3_SHAPEMORPHING}장에서 쓴 종류의 로그 질량 격자 손실과 같은 기울기 노름을 주며, 이로써 정리 가중치 $w_c$의 단위도 정해진다. 제어나 부피에 대한 정칙화 항은 없다.

\subsection{물리 항이 묻는 것}
\emph{수송} $\mathcal{L}_{\mathrm{OT}}$은 손실 격자 위에서 몸체 질량과 목표 질량 사이의 편향 제거 Sinkhorn
발산~\citep{feydy2019sinkhorn}이다. 이 선택은 모핑에 중요하다. \ref{CH3_SHAPEMORPHING}장의 노드 질량 잔차는 기울기가 표면에서
밀도가 뛰는 곳에만 있으므로 단지 엉뚱한 곳에 있을 뿐인 채워진 영역을 보지 못하지만, 수송의 기울기는 몸체 전체에 걸친
수송장이며 위상에 무관하다. 수송은 손실 셀 안의 위치를 흐리므로, \emph{표면 근접} $\mathcal{L}_{\mathrm{prox}}$이 목표의 모든
바깥 점에서 그곳에 몸체 물질이 있는지 묻는다. 두 안정 항은 정지한 끝을 요구한다. \emph{끝 드리프트}
$\mathcal{L}_{\mathrm{drift}}$는 해제된 끝에 남은 속도에, \emph{해제 운동} $\mathcal{L}_{\mathrm{rel}}$은 해제 절반 동안의 모든
운동에 비용을 매긴다. 해제 운동이 없으면 몸체가 아직 움직이는 채로 실행이 멈춘다. \emph{정리} 묶음은 선택 사항이다. 두
국소 항, 식~\eqref{ch4:eq_spray_ko}와 \eqref{ch4:eq_near_ko}는 수송이 손실 셀 아래에서 해상하지 못하는 것, 곧 몸체에서
흩어져 나간 입자와 목표 표면 바로 바깥의 얇은 띠에 있는 입자만 정리한다. 우리는 늘 $w_c=0.2$를 쓰며, $w_c=0$이면 이 묶음이
빠진다. 목적 함수를 정할 때의 하나씩 빼기 실험에서 흩뿌림 당김을 빼면 형상은 그대로였지만 고립된 입자가 늘고 실행이
길어졌으며, 근접 띠를 빼면 입자 40k에서 덮이지 않은 얇은 목표 점이 늘었다(부록~\ref{ch4:sec_cleanup_ko}).

\subsection{외피 위의 렌더 항}
\label{ch4:sec_exterior_ko}
렌더 항은 입자 구름이 아니라 몸체의 표면을 읽는다. 몸체의 \emph{외피}는 입자에 대한 Zhu와 Bridson의 장~\citep{zhu2005sand}의
영점 집합이며, 입자를 따라가는 표면 원판이 덮는다. \emph{실루엣} 항은 18개의 직교 시점에서 원판의 덮임을 목표의 것과
비교하되 목표 안쪽의 덮임 부족을 바깥쪽의 초과보다 더 크게 벌하고, \emph{셰이딩} 항은 헤드라이트로 비춘 원판의 영상을
비교한다. 목표 그림은 목표 샘플의 외피에서 같은 연산자로 그리되, 메시의 독립 샘플 여덟 개에 대해 평균하므로 기준에 어느 한
샘플의 잡음도 들어 있지 않다(부록~\ref{app:ch4_render_ko}). 렌더 항을 입자 구름이 아니라 외피에서 읽기 때문에 렌더 항이
관찰자가 보는 표면에 작용한다. 구름에서 읽으면 구름 자신의 실루엣은 나아지지만 표시되는 표면은 쌍둥이의 자리에 남는다
(\ref{ch4:sec_ablation_ko}절). 이 장의 그림은 같은 외피를 표면 원판으로 스튜디오 룩에 그리며, 모든 그림은 확정된 시뮬레이션
상태이다.

% ============================================================================
\section{모핑 방법}
\label{ch4:sec_method_ko}
% ============================================================================

몸체, 제어, 목적 함수가 주어지면, 모핑은 \ref{CH3_SHAPEMORPHING}장처럼 짧은 시뮬레이션 네트워크들의 연속에 대한
최적화이다. 여기서 각 네트워크는 안정화 창이며, 창은 몸체를 풀어 놓고 평가한 뒤에만 확정된다.

\subsection{안정화 창}
한 창은 제어가 있는 $T$개의 구동 단계, 그다음 제어를 0으로 두고 물리만 계속 돌리는 $T$개의 해제 단계를 시뮬레이션한다.
목적 함수의 모든 항은 해제 후의 끝 상태와 해제 중의 속도에서 계산한다. 따라서 되튐, 진동, 그 밖에 제어가 밀고 있는 동안만
유지되는 것은 모두 제어에 직접 비용이 된다. 최적화기는 창마다 정해진 횟수만큼 반복하고, 그다음 창을 확정하며, 다음 창은
확정된 상태에서 시작한다(알고리즘~\ref{ch4:alg_window_ko}).

\subsection{기울기와 렌더 가중치}
\label{ch4:sec_lambda_ko}
물리 항, 정리 항, 렌더 항은 같은 $2T$단계 수반을 통해 모든 $\delta\mathbf{F}^n_{c}$와 $u$까지 따로 미분한다.
$\mathbf{g}_p$를 물리 핵심의 기울기, $\mathbf{g}_r$을 렌더 항의 기울기라 하자. 렌더 기울기는 둘이 충돌할 때 물리 기울기에
반대되는 성분을 잃고~\citep{yu2020pcgrad},
\begin{equation}
\tilde{\mathbf{g}}_r=\mathbf{g}_r-\frac{\min(0,\langle\mathbf{g}_p,\mathbf{g}_r\rangle)}{\lVert\mathbf{g}_p\rVert^{2}}\,\mathbf{g}_p ,
\end{equation}
렌더 가중치는 각 창의 첫 기울기에서
\begin{equation}
\label{ch4:eq_lambda_ko}
\lambda\,\lVert\tilde{\mathbf{g}}_r\rVert=\tfrac12\,\lVert\mathbf{g}_p\rVert
\end{equation}
로 정하므로, 렌더 항은 모든 창에서 첫 기울기의 3분의 1을 차지한다. 단계 방향은
$\mathbf{g}=\mathbf{g}_p+\lambda\tilde{\mathbf{g}}_r+\mathbf{g}_{\mathrm{clean}}$이다.

이 규칙은 한 단계에서 렌더가 차지하는 \emph{몫}은 일정하게 유지하지만, 그 \emph{권한}은 유지하지 않는다. 몸체가 도착하는
동안 물리 기울기는 여러 자릿수 떨어지고, $\lambda$도 함께 떨어진다. 그래서 렌더 항은 몸체가 오는 동안에는 강하게, 도착한
뒤에는 약하게만 작용한다. 이 결과를 \ref{ch4:sec_render_effect_ko}절에서 재고, \ref{ch4:sec_costs_ko}절에서 다시 다룬다.

\subsection{두 번째 방식: 렌더가 가중치를 유지한다}
\label{ch4:sec_held_ko}
식~\eqref{ch4:eq_lambda_ko}의 규칙은 렌더의 가중치를 물리 잔차에 묶는다. 300k bunny에서 $\lambda$는 첫 창의 0.32에서 몸체가
도착한 뒤 0.003 아래로 떨어지고, 바깥층 당김 중 렌더의 몫은 5분의 1 근처에 머문다. 그래서 렌더만 정할 수 있는 것, 곧 어느
물질이 어디에 끝나는지와 3D 목표가 해상하는 것보다 미세한 부조가 남았을 바로 그때 렌더에는 권한이 남아 있지 않다. 모든
물리 항은 입자 위치의 집합만 읽으므로 물리 목적 함수는 둘 다 보지 못한다.

두 번째 운용 방식에서는 렌더가 첫 창의 가중치를 유지한다. $\lambda$를 식~\eqref{ch4:eq_lambda_ko}로 한 번 보정한 뒤 묶거나,
그림에만 있는 내용을 담는 항(예: 색 항)의 가중치만 묶고 나머지 렌더 항은 규칙을 따르게 한다. 이 방식의 목적은 그림만 담고
있는 내용, 곧 목표 그림에만 있는 부조와 겉모습에 따른 물질 배치이다. 도착 뒤에는 고정 가중치의 정리 항도 렌더에 맞서 미는
힘의 하나이므로 $w_c=0$이 이 방식에 자연스러운 짝이지만, \ref{ch4:sec_write_ko}절의 실행들은 정리 항을 유지한다. 대가도
있다. 묶인 가중치는 몸체가 도착해도 사라지지 않으므로 실행이 멈출 때 몸체가 아직 움직이고, 가장 촘촘한 형상 척도가 조금
느슨해진다. 렌더가 정할 것이 있는 곳에서는 권한을 유지하고 몸체는 정지하게 하는 가중치는 진행 중인 작업이다
(\ref{ch4:sec_costs_ko}절).

\subsection{단계, 확정, 정지}
단계는 \ref{CH3_SHAPEMORPHING}장처럼 되돌아가는 선 탐색을 붙인 Adam~\citep{kingma2017adammethodstochasticoptimization}이고,
각 시도 뒤 $\delta\mathbf{F}_c$에 상한을 두고 $u$는 크기를 제한하고 강체 성분을 없앤다. $u$의 모멘트는 다음 창으로 이어지므로,
창의 첫 단계가 앞 창의 단계를 뒤집지 않는다. 창의 결과는 소성 동화와 함께 확정되고, 목적 함수를 크게 올리면 거절되며,
목적 함수가 더 나아지지 않으면 실행이 멈춘다(부록~\ref{app:ch4_method_ko}).

\begin{algorithm}[t]
  \caption{안정화 창 모핑}\label{ch4:alg_window_ko}
  \begin{algorithmic}[1]
    \Require 원천과 목표 샘플, 구동 $T$단계와 해제 $T$단계, 창마다 반복 $i_{\max}$회
    \State $\mathbf{s}\gets$ 정지한 원천 몸체;\ $\delta\mathbf{F}^{\mathrm{prev}}_c\gets\mathbf{0}$
    \While{정지 규칙이 작동하지 않음}
      \State 0에서 시작하는 것보다 나으면 $\delta\mathbf{F}_c\gets\tfrac12\delta\mathbf{F}^{\mathrm{prev}}_c$, 아니면 $\mathbf{0}$;\ $u\gets 0$
      \For{$i=1,\dots,i_{\max}$}
        \State $\mathbf{s}$에서 구동 $T$단계($\delta\mathbf{F}_c$, $u$ 작용)와 해제 $T$단계를 시뮬레이션
        \State 해제된 끝에서 $\mathcal{L}_{\mathrm{phys}}$, $\mathcal{L}_{\mathrm{clean}}$, $\mathcal{L}_{\mathrm{render}}$ 계산
        \State $\mathbf{g}_p,\mathbf{g}_{\mathrm{clean}},\mathbf{g}_r\gets$ $2T$단계를 거치는 수반
        \State $\tilde{\mathbf{g}}_r\gets$ PCGrad$(\mathbf{g}_r,\mathbf{g}_p)$;\ $i=1$일 때 $\lambda\gets\tfrac12\lVert\mathbf{g}_p\rVert/\lVert\tilde{\mathbf{g}}_r\rVert$
        \State $\mathbf{g}_p+\lambda\tilde{\mathbf{g}}_r+\mathbf{g}_{\mathrm{clean}}$에 대해 선 탐색을 붙인 Adam 단계
        \State $\delta\mathbf{F}_c$ 상한;\ $u$ 크기 제한과 강체 모드 제거
      \EndFor
      \State 받아들인 후보 재현;\ 탄성 늘어남의 절반을 $\mathbf{F}^{\mathrm{P}}$에 동화
      \If{기준이 $5\,\%$ 넘게 올랐음}
        \State 거절:\ $\mathbf{s}$ 유지
      \Else
        \State 확정:\ $\mathbf{s}\gets$ 해제된 끝 상태;\ $\delta\mathbf{F}^{\mathrm{prev}}_c\gets\delta\mathbf{F}_c$
      \EndIf
    \EndWhile
    \State \Return 가장 좋은 창까지의 궤적
  \end{algorithmic}
\end{algorithm}

\subsection{몸체가 보존하는 것}
\label{ch4:sec_conservation_ko}
질량은 구성상 정확하다. 응력 제어는 선운동량도 각운동량도 바꿀 수 없고, 최소 간격의 밀어냄은 쌍별 중심력이며, 이완과 $u$는
바깥층의 강체 모드를 갖지 않는다. 예외는 약한 주변 항력 하나로, 각운동량과 일관되지 않으므로 회전은 운동학적 양으로
보고한다. 300k 모핑 전체에서 질량 중심은 $10^{-4}$ 피치 미만 움직이고, 몸체의 물리적 부피는 목표의 0.5\,\% 안에서 끝난다.

\subsection{장면 속의 몸체}
세계에 고정된 목적 함수 하나로는 몸체가 날거나, 돌거나, 착지할 수 없다. 여러 몸체, 접촉, 중력이 있는 장면에서는 각 몸체가
자기 격자장에 질량을 쌓아 달라붙는 대신 운동량을 보존하는 접촉 규칙으로 만나고, 각 몸체의 목적 함수를 자기 형상 정합
좌표계에서 읽어 비행과 구르기에 비용이 들지 않게 한다. 모든 창을 확정하며, 목표는 수송을 따라 움직여 몸체들이 반쯤 만들어진
상태로 만난다. 한 몸체 경로는 이 확장으로 바뀌지 않는다(부록~\ref{ch4:sec_scene_ko}).

% ============================================================================
\section{예제}
\label{ch4:sec_examples_ko}
% ============================================================================

\subsection{설정}
\label{ch4:sec_protocol_ko}
주 평가의 모든 모핑은 구에서 시작하여 $N=300{,}000$에서 bunny, dragon, teapot, bob, armadillo, homer 중 하나를 목표로 하며,
그중 셋은 손잡이가 있다. 19개 메시 갤러리는 $N=40{,}000$에서 돌리고, 목표와 형상 쌍 20개를 더 300k에서 돌렸다. 비교하는
갈래는 담고 있는 것으로 이름 붙인다. \emph{전체 방법}($\delta\mathbf{F}_c$ + layer $u$ + render loss), \emph{그 $\lambda=0$
쌍둥이}($\delta\mathbf{F}_c$ + layer $u$), \emph{$\delta\mathbf{F}_c$만}, 그리고 \emph{\ref{CH3_SHAPEMORPHING}장의 F-제어}이다.

어떤 결과도 그것이 최적화한 샘플에 대고 평가하지 않는다. 각 끝을 학습이 보지 않는 시점에서 정확한 목표 메시에 대고 표준
실루엣·표면·영상 척도로 읽고, 입자 수준에서 메시의 독립 샘플에 대고 읽으며(두 번째 샘플이 \emph{샘플링 바닥}, 즉 완벽한
배치를 정한다), 렌더 항 자신의 연산자로도 읽는다. 마지막 것은 렌더가 어디서 작용하는지 찾아 주지만, 아래에서 보이듯
바깥층의 배열에도 보상을 준다. 부록~\ref{app:ch4_measures_ko}에 각 척도와, 어떤 척도도 계산하기 전에 정한 ``렌더를 본다''의
규칙을 적었다.

\begin{figure*}[t]
    \centering
    \includegraphics[width=\textwidth]{figs/ch4/fig_morph_sequence_v2.jpg}
    \caption{입자 300k에서 구에서 다섯 목표로. \textmd{모든 줄을 같은 시각들과 끝에서 보인다. 몸체는 한 덩어리로 솟고 첫
    1--2초에 특징이 그로부터 자라 나오며, bob의 고리는 조금 뒤에, dragon의 가장 좁은 틈은 마지막에 열린다. 입자 10개 미만을
    담은 표시 표면의 따로 떨어진 조각은 그리지 않는다(부록~\ref{app:ch4_render_ko}; 패널 글자는 영어판과 공유).}}
    \label{ch4:fig_morph_ko}
\end{figure*}

\subsection{구에서 다섯 목표로}
그림~\ref{ch4:fig_morph_ko}에 다섯 목표를 보였다. 몸체는 한 덩어리로 솟고, 첫 1--2초 안에 특징이 그로부터 자라 나온다.
bunny의 귀와 armadillo의 팔다리는 1.5\,s에 자리를 잡고, bob의 고리는 조금 뒤에 열리며, dragon의 뿔과 발톱, 똬리 사이의
좁은 틈은 마지막에 온다. 입자 이동의 대부분은 이 첫 몇 초에 일어나고, 그 뒤에는 결과 프레임의 대부분을 채우는 길고 느린
꼬리가 온다. 끝은 입자를 완벽하게 배치했을 때 얻을 값에 가깝다. 입자 실루엣은 샘플링 바닥 바로 아래에, chamfer 거리는
바닥보다 아래에 있는데, 끝이 어떤 무작위 샘플보다 고르기 때문이다. 몸체는 정지한 채 끝나고, 어느 실행도 심각한 조각을
떨어뜨리지 않는다. 쌍둥이와 비교하면 전체 방법은 모든 메시와 시드에서 정확한 목표에 더 가깝게 끝나며, 그 작은 차이는 얇은
부분의 테두리에 있다(그림~\ref{ch4:fig_render_effect_ko}). \ref{ch4:sec_write_ko}절은 렌더가 도착 뒤에도 권한을 유지할 때
무엇을 쓰는지 보이고, \ref{ch4:sec_render_effect_ko}절은 두 방식에서 렌더가 어떻게 작용하는지 따라간다.

\begin{figure*}[tbp]
    \centering
    \includegraphics[width=0.9\textwidth]{figs/ch4/fig_render_effect_v2.png}
    \caption{300k 모핑의 끝에서 렌더 손실이 바꾸는 것. \textmd{(a, b) 메시와 시드마다 한 줄이며, 학습이 보지 않는 시점에서
    정확한 목표 메시에 대고 읽는다. 회색은 $\lambda=0$ 쌍둥이($\delta\mathbf{F}_c$ + layer $u$), 주황은 전체 방법, 검은 눈금은
    목표의 완벽한 입자 샘플이다. 윤곽 오차는 22쌍 모두에서 중앙값 16\,\% 낮고, 반 피치에서의 표면 오차는 22쌍 모두에서 중앙값
    36\,\% 낮다. (c) 차이는 얇은 부분의 테두리에 있다. bunny의 귀와 dragon의 꼬리 끝의 끝 실루엣(시드 97, 홀드아웃 시점)을
    목표의 윤곽과 겹쳤다(패널 글자는 영어판과 공유).}}
    \label{ch4:fig_render_effect_ko}
\end{figure*}

\subsection{렌더가 쓸 수 있는 것}
\label{ch4:sec_write_ko}
그림~\ref{ch4:fig_render_effect_ko}는 \ref{ch4:sec_lambda_ko}절의 기본 규칙 아래의 렌더이다. 그 가중치는 물리 기울기를
따르므로, 몸체가 도착하면 렌더는 테두리를 다듬기만 한다. 이것은 가중치의 성질이지, 렌더가 닿을 수 있는 범위의 성질이
아니다. 도착 뒤에도 가중치를 묶어 두면(\ref{ch4:sec_held_ko}절) 렌더는 물리 목적 함수가 할 수 없는 일을 한다. 그림만 담고
있는 내용을 쓴다.

\textit{그림에만 있는 부조.} 3D 목표는 매끈한 판이고, 앞면의 SIGGRAPH와 뒷면의 VANDERBILT라는 글자는 렌더의 목표 그림에만
있다(그림~\ref{ch4:fig_letters_ko}). 바깥층 오프셋 $u$는 있지만 렌더 손실이 없는 쌍둥이는 매끈하게 끝난다. 렌더 가중치를
묶으면 전체 방법은 두 단어를 입자에 새기며, 글자가 있는 메시 자체를 300k로 샘플링한 것이 담는 부조의 58\,\%에 이른다. 이는
보이는 부조가 $u$만으로 생길 수 있는지에도 답한다. $u$는 두 갈래 모두에서 작용하고, 그림을 보는 갈래만 글자를 쓴다. 규칙
아래에서는 판이 만들어질 즈음 렌더의 가중치가 이미 사라져, 40k 판에서 글자는 쌍둥이 수준에 머문다. 비용은
\ref{ch4:sec_held_ko}절에서 말한 그대로이다. 실행이 멈출 때 몸체는 아직 움직이고 있고, 판의 양 끝이 조금 덜 채워진다.
%% SWAP-IN: 300k 물결(ripple) 부조 결과(그림과 숫자 하나)가 나오면 여기에 넣는다.
%% SWAP-IN: u-off 답("dF_c만", 판에서 렌더 가중치 고정)이 나오면 여기에 넣고, "u는 두 갈래 모두에서 ..." 문장을
% 직접 측정으로 바꾼다.

\begin{figure*}[tbp]
    \centering
    \includegraphics[width=\textwidth]{figs/ch4/fig_render_letters.jpg}
    \caption{목표 그림에만 있는 글자. \textmd{3D 목표가 매끈한 300k 판. 렌더 가중치를 묶으면 전체 방법은 SIGGRAPH(앞면)와
    VANDERBILT(뒷면)를 입자에 새기고, 바깥층 오프셋 $u$를 가진 쌍둥이는 매끈하게 끝난다(패널 글자는 영어판과 공유).}}
    \label{ch4:fig_letters_ko}
\end{figure*}

\textit{겉모습에 따른 배치.} 각 입자는 원천에서 정해져 함께 움직이는 색을 가지며, 목표 그림은 색칠되어 있다. 물리 목적
함수는 이 색칠을 증명 가능하게 보지 못한다. 쌍둥이의 궤적은 어떤 색칠에서도 비트 단위로 같다. 그림~\ref{ch4:fig_colour_ko}의
장면에서 원천은 위쪽 절반이 빨갛고, 그림은 bunny의 머리 쪽 절반을 빨갛게 칠한다. 쌍둥이는 수송이 만든 위아래 나눔을 유지하고,
규칙 아래의 전체 방법도 배치를 쌍둥이와 거의 같은 곳에 둔다. 색 항의 가중치를 묶으면 전체 방법은 그림이 칠한 곳에 색이
놓이도록 물질을 옮기고, 배치 오차의 약 70\,\%를 되찾는다. 이 실행도 색이 아직 움직이는 채로 끝나며, 가장 촘촘한 표면 척도가
조금 느슨해진다. 이 절의 측정값은 부록~\ref{app:ch4_held_ko}에 있다.
%% SWAP-IN: 300k 색 장면(그림과 숫자 하나)이 나오면 여기에 넣는다.

\begin{figure*}[tbp]
    \centering
    \includegraphics[width=\textwidth]{figs/ch4/fig_render_colour.jpg}
    \caption{겉모습에 따라 물질을 배치하기. \textmd{물리 목적 함수는 어느 입자가 어떤 색을 가졌는지 보지 못한다. 색 항의
    가중치를 묶으면 전체 방법은 빨간 물질을 목표 그림이 칠한 머리 쪽 절반으로 옮기고, 쌍둥이는 수송의 위아래 나눔을 유지한다
    (bunny, 입자 40k, 앞과 뒤; 패널 글자는 영어판과 공유).}}
    \label{ch4:fig_colour_ko}
\end{figure*}

\subsection{렌더 항이 하는 일}
\label{ch4:sec_render_effect_ko}
규칙 아래에서 렌더 항을 쌍둥이와 비교하며 실행 전체에 걸쳐 따라간다. 도착하는 동안, 도착한 뒤, 그리고 끝의 세 국면이며, 그다음
가중치를 묶었을 때 렌더가 작용하는 곳을 비교한다.

\textit{도착하는 동안, 렌더는 물질의 방향을 바꾼다.} 첫 1초 동안 전체 방법의 표면은 쌍둥이의 표면에서, 같은 갈래를 다른
시드로 두 번 돌린 실행끼리보다 훨씬 크게 벌어진다. dragon에서 그 간격은 거의 한 피치에 이른다. 그 1초 동안 렌더가 나르는
몫은 표면 운동의 작은 부분에 불과하지만, 그 1초가 모핑이 하는 운동의 대부분을 담는다. 그림~\ref{ch4:fig_rgv_ko}는 렌더가
작용하는 곳을 보인다. 귀, 얼굴, 발이며, 초반에는 대체로 물리와 같은 쪽으로, 후반에는 더 자주 반대쪽으로 당기고, 그곳에서
투영이 반대 성분을 제거한다.

\textit{도착한 뒤, 렌더는 다듬기만 한다.} 몸체가 도착하면 물리 기울기는 한 자릿수 넘게 떨어져 있고, \ref{ch4:sec_lambda_ko}절의
규칙에 따라 렌더 가중치와 단계도 함께 떨어진다. 두 갈래의 표면 차이는 시드 간 잡음보다 크지 않은 수준으로 돌아온다. 렌더
항이 여전히 하는 일은 실루엣 테두리와 얇은 부분에서 물리 자신의 재맞춤에 맞서 피치 이하의 오프셋을 붙잡는 것이다. 인과
실험이 이를 보인다(그림~\ref{ch4:fig_switch_ko}). 6\,s에 끄면 전체 방법은 약 1초 안에 앞섬을 잃고, 6\,s에 켜면 쌍둥이는 3분의
1초 안에 앞섬을 만들어 처음부터 렌더가 있던 실행과 같은 수준으로 끝난다. 렌더가 요구하는 것의 대부분은 MPM 셀 아래에 있고,
그것에 닿는 것은 $\delta\mathbf{F}_c$가 아니라 $u$이다.

\textit{끝에서 이득은 작지만 일관된다.} 표~\ref{ch4:tab_d157_ko}는 학습이 보지 않는 시점에서 정확한 목표 메시에 대고 끝을
읽는다. 300k에서는 하나를 뺀 모든 표준 척도가 모든 메시와 거의 모든 실행 쌍에서 나아지며, 예외인 하나는 포화되어 있다. 개선은
실재하지만 절댓값으로는 작다. Boundary IoU 오차는 0.068에서 0.054로, 완벽한 샘플의 0.052 가까이 간다
(그림~\ref{ch4:fig_render_effect_ko}). 모든 학습 시점에서 멀리 떨어진 시점에서도 이득이 같으므로 학습이 보는 그림에만 맞춘
것이 아니며, 표~\ref{ch4:tab_b1_ko}에 메시별로 보였다. 40k에서는 같은 척도가 같은 방향으로, 경계 척도에서는 약 절반,
나머지에서는 그보다 적게 움직인다. 입자 수준의 척도는 이 이득을 보지 못한다. 입자 실루엣, chamfer 거리, 얇은 부분 덮임, 질량
중심, 정지는 쌍둥이의 범위 안에 있는데(그림~\ref{ch4:fig_render_end_ko}), 이진 입자 실루엣은 한 픽셀에 거의 두 피치를 담는
반면 렌더는 외피를 피치의 일부만큼 바꾸기 때문이다(그림~\ref{ch4:fig_iou_ko}).

\begin{table}[t]
\centering
\footnotesize
\setlength{\tabcolsep}{4pt}
\caption{정확한 목표 메시에 대고 읽은 전체 방법 대 그 $\lambda=0$ 쌍둥이. \textmd{$r$은 쌍둥이 오차 중 전체 방법이 없앤 비율의
중앙값(PSNR은 dB 이득)이고, $^\ast$는 미리 정한 규칙상 렌더를 보는 척도이다.}}
\label{ch4:tab_d157_ko}
\begin{tabular}{lcccc}
\toprule
 & \multicolumn{2}{c}{300k (메시 6, 쌍 22)} & \multicolumn{2}{c}{40k (메시 19, 쌍 57)} \\
\cmidrule(lr){2-3}\cmidrule(lr){4-5}
척도 & $r$ & 쌍 / 메시 & $r$ & 쌍 / 메시 \\
\midrule
\multicolumn{5}{l}{\textit{실루엣, 홀드아웃 시점 (1024\,px)}}\\
Boundary IoU (대각선의 2\,\%)          & $+16.3\,\%^\ast$ & 22 / 6 & $+7.1\,\%^\ast$ & 52 / 17 \\
경계 F-score (DAVIS, 12\,px)            & $+0.9\,\%$       & 12 / 3 & $+12.5\,\%$      & 31 / 11 \\
윤곽 Chamfer                            & $+20.2\,\%^\ast$ & 21 / 6 & $+8.4\,\%^\ast$  & 54 / 18 \\
윤곽 HD95                               & $+19.6\,\%^\ast$ & 21 / 6 & $+8.1\,\%^\ast$  & 50 / 18 \\
\multicolumn{5}{l}{\textit{3D, 보이는 외피}}\\
F-score, 0.5 피치                       & $+36.4\,\%^\ast$ & 22 / 6 & $+8.0\,\%^\ast$  & 46 / 16 \\
F-score, 1 피치                         & $+14.4\,\%^\ast$ & 22 / 6 & $+3.5\,\%$       & 39 / 15 \\
법선 일관성                             & $+9.2\,\%^\ast$  & 21 / 6 & $+2.8\,\%$       & 42 / 15 \\
Chamfer-L1                              & $+12.9\,\%^\ast$ & 22 / 6 & $+4.3\,\%$       & 46 / 16 \\
\multicolumn{5}{l}{\textit{영상, 홀드아웃 시점 (512\,px)}}\\
셰이딩 PSNR                             & $+0.73$\,dB$^\ast$ & 22 / 6 & $+0.30$\,dB$^\ast$ & 53 / 18 \\
셰이딩 SSIM                             & $+10.7\,\%^\ast$ & 22 / 6 & $+3.0\,\%$       & 52 / 18 \\
셰이딩 LPIPS                            & $+7.0\,\%^\ast$  & 22 / 6 & $+2.5\,\%$       & 47 / 17 \\
법선 지도 PSNR                          & $+0.96$\,dB$^\ast$ & 22 / 6 & $+0.36$\,dB$^\ast$ & 53 / 18 \\
법선 지도 SSIM                          & $+9.0\,\%^\ast$  & 22 / 6 & $+2.7\,\%$       & 51 / 17 \\
법선 지도 LPIPS                         & $+10.0\,\%^\ast$ & 22 / 6 & $+3.8\,\%$       & 51 / 17 \\
\midrule
\multicolumn{5}{l}{\textit{비교용}}\\
입자 실루엣 IoU (128\,px)               & $+2.3\,\%$       & 14 / 4 & $+4.5\,\%$       & 48 / 19 \\
원래 메시에 대한 마스크 IoU             & $+5.2\,\%$       & 22 / 6 & $-0.9\,\%$       & 19 / 5 \\
\bottomrule
\end{tabular}
\end{table}

\textit{렌더 자신의 항에서 앞선 몫의 일부는 형상이 아니라 배열이다.} 300k 실행의 끝에서 전체 방법의 렌더 항은 진짜 메시의
어떤 균일 샘플보다도 낮다. 정확한 형상을 균일하게 샘플링한 몸체는 그렇게 할 수 없다. 그 바닥 아래로 내려가는 것은 어떤
무작위 샘플보다 고르게 배열된 바깥층, 따라서 더 매끈한 외피이며, 렌더 항에서 전체 방법이 쌍둥이보다 앞선 몫 중 bunny에서는
최소 5분의 2, dragon에서는 최소 5분의 1이 이런 종류이다. 그래서 형상은 정확한 메시로만 읽고, 그림~\ref{ch4:fig_render_end_ko}와
\ref{ch4:fig_switch_ko}처럼 렌더 항 자신의 연산자는 렌더의 작용을 찾는 데만 쓴다.

요컨대 규칙 아래에서 렌더 항은 도착을 바꾸고 끝을 다듬지만, 도착한 뒤에는 물리 목적 함수가 열어 둔 것을 아무것도 정하지
않는다. 가중치를 묶으면 바뀌는 것이 바로 이것이다(\ref{ch4:sec_write_ko}절). 작은 비용도 있다. 일부 몸체를 더 변형시키고,
각운동량을 지키지 않는 주변 항력은 더 움직이는 몸체에 더 작용한다. 이 절의 측정값은 부록~\ref{app:ch4_results_ko}에 있다.

\begin{figure*}[t]
    \centering
    \includegraphics[width=\textwidth]{figs/ch4/fig_render_switch.png}
    \caption{6\,s에서 렌더 항을 전환. \textmd{끄면 전체 방법은 약 1초 안에 앞섬을 잃고, 켜면 쌍둥이는 3분의 1초 안에 앞섬을
    만든다.}}
    \label{ch4:fig_switch_ko}
\end{figure*}

\begin{figure*}[p]
    \centering
    \includegraphics[width=\textwidth]{figs/ch4/fig_render_gradient_early.jpg}\\[4pt]
    \includegraphics[width=\textwidth]{figs/ch4/fig_render_gradient_late.jpg}
    \caption{변형 중인 300k bunny 위의 렌더 기울기, 초반(위)과 후반(아래). \textmd{각 블록은 입자별 당김 중 렌더의 몫, 그 세기,
    물리와의 일치, 그리고 렌더가 $\delta\mathbf{F}_c$에서 바꾼 것을 보인다(패널 글자는 영어판과 공유).}}
    \label{ch4:fig_rgv_ko}
\end{figure*}

\textit{가중치를 묶었을 때 렌더가 작용하는 곳.} 그림~\ref{ch4:fig_held_share_ko}는 그림~\ref{ch4:fig_rgv_ko}의 몫 지도를 300k
bunny에서 규칙 아래와 렌더 가중치를 묶은 경우로 되풀이한다. 규칙 아래에서 바깥층 당김 중 렌더의 몫은 도착 뒤 0.14--0.23에
머물고, 묶으면 0.78--0.90이 되어, 그 뒤로는 물리가 아니라 렌더가 바깥층이 할 일을 정한다. 렌더가 \ref{ch4:sec_write_ko}절의
내용을 쓰는 방식이 이것이고, 그 대가도 여기서 온다. 끝의 운동 에너지가 규칙 아래보다 세 자릿수 높고, 덮이지 않은 얇은 목표
점이 조금 늘어난다(부록~\ref{app:ch4_held_ko}).
%% SWAP-IN: "lambda 고정 + 손실 셀에서 읽는 배치 항" 갈래가 나오면 여기에 넣는다. 글자를 유지하면서 몸체를
% 정지시키면 위의 비용 문장도 바꾼다.

\begin{figure}[tbp]
    \centering
    \includegraphics[width=0.62\linewidth]{figs/ch4/fig_render_held_share.jpg}
    \caption{가중치를 묶었을 때 렌더가 작용하는 곳. \textmd{300k bunny에서 각 입자의 당김 중 렌더의 몫, 회색(물리만)에서
    주황(렌더만)까지. 규칙 아래에서는 도착 뒤 옅어지고, 묶으면 바깥층을 지배한다(패널 글자는 영어판과 공유).}}
    \label{ch4:fig_held_share_ko}
\end{figure}

\subsection{40k 갤러리}
\label{ch4:sec_gallery40_ko}
$N=40{,}000$의 19개 메시 갤러리에서 렌더는 300k와 같은 방향으로, 약 절반 크기로 작용한다. 입자 실루엣은 하나를 뺀 모든
메시에서 쌍둥이보다 작은 차이로 앞서고, 정확한 메시에 대고 보면 윤곽은 19개 모두에서 나아진다. 한 판독은 반대로 간다. 원래
메시에 대한 실루엣 면적이 조금 나쁜데, 외피가 조금 바깥으로 기울기 때문이다. 이 기울어짐은 덮임 부족과 초과를 비대칭으로
가중하는 실루엣 항에서 오며, 300k와 달리 40k에서는 안쪽으로 당기는 힘이 이를 상쇄하지 않는다. 측정과 다음 시험은
부록~\ref{app:ch4_gallery_ko}에 있다.

\subsection{300k 모핑 모음}
\label{ch4:sec_showcase_ko}
여섯 주 목표 밖에서도 방법을 보기 위해, 방법을 바꾸지 않고 구에서 갤러리의 다른 13개 메시와 genus 6인 Happy Buddha로, 그리고
여섯 형상 쌍 사이에서 돌렸다(그림~\ref{ch4:fig_showcase_ko}). 하나를 뺀 모든 실행이 주 목표 여섯과 거의 같은 정도로 목표에
도달하고, 정지하며, 심각한 조각을 떨어뜨리지 않는다. 위상은 목표가 요구하는 대로 바뀐다. teapot에서 bob으로 갈 때 teapot의
손잡이는 첫 1초 안에 닫히고 bob의 고리는 나중에 열려, 몸체는 목표의 구멍 하나를 가지고 끝난다. 구에서 시작한 Happy Buddha에서는
구멍들이 생긴다. 쌍둥이를 돌린 경우 렌더의 앞섬은 주 목표 여섯에서와 같은 방향, 같은 크기이다.

세 예제는 방법이 아직 치르는 비용을 보인다. 글자 C에서는 몸체가 아직 흐르는 채로 실행이 일찍 멈춘다. 정지 규칙이 정지가 아니라
목적 함수의 진전을 읽기 때문이다. spot과 cow처럼 방향이 다른 두 형상 사이에서는 수송이 머리를 등 위로 넘긴다. 원천을 먼저
돌려 두면 같은 모핑이 곧바르고 두 배 빨라진다. 그리고 armadillo처럼 덜 조밀한 원천은 몸체에 비해 거친 셀에서 돌며, homer로
가는 도중 구멍 난 얇은 판을 거친다.

\begin{figure*}[p]
    \centering
    \includegraphics[width=\textwidth]{figs/ch4/fig_showcase_montage_v2_interim24.jpg}
    \caption{300k 모핑 24개, 각각 시작·중간·끝. \textmd{마지막 패널은 방법 밖의 옵션인 원천 부피로 정한 MPM 셀을 써서 바로 앞
    패널과 비교한다(패널 글자는 영어판과 공유).}}
    \label{ch4:fig_showcase_ko}
\end{figure*}

\subsection{재질과 외력}
시험한 모든 재질이 목표에 도달한다. 무른 고체는 더 천천히, 단단한 고체는 더 빨리 그러나 덜 정지한 채 도착하며, 렌더 항의
앞섬은 각 재질에서 유지된다. 중력 아래에서 한 몸체 창은 일찍 멈춘다. 떨어지는 몸체가 세계에 고정된 목적 함수를 올리기
때문이다. 목표를 바닥 위에 세워 둔 장면 경로에서는 bunny와 dragon 모두 모든 시드에서 목표에 도달하고, 목표의 높이에 똑바로
서며, 수평 운동량을 지킨다. 다만 정지하지는 않는다. 매 확정의 동화가 무게로 인한 처짐의 일부를 정지 형상으로 만들고 제어가
이를 다시 들어 올리기 때문이다. 다음 정의는 각 창의 변형 중 제어의 몫만 동화한다(부록~\ref{app:ch4_gravity_ko}).

\subsection{제거 실험}
\label{ch4:sec_ablation_ko}
300k에서 한 번에 한 부분씩 끄면(표~\ref{ch4:tab_ablation_ko}) 세 가지 판독이 나온다. 첫째, $u$가 주된 형상 지렛대이다. 없으면
모든 표면 척도가 몇 배 나빠지고 몸체는 끝까지 정지하지 않지만, 초기 떠다니는 입자는 거의 사라진다. $u$가 그 원천이기도 하기
때문이다. 둘째, 렌더를 외피에서 읽으면 표면의 모든 척도에서 입자 구름에서 읽는 것보다 낫고, 구름 읽기는 자신이 최적화하는
입자 척도에서 낫다. 각 읽기는 자신이 평가하는 표현에서 이기며, 외피가 관찰자가 보는 표현이다. 셋째, 부피 보존 몸체, 최소
간격, 늦은 운동에 맞선 선택들은 각자 겨냥한 것에서 제 몫을 한다.

\subsection{응력 제어 대 힘 제어}
\label{ch4:sec_force_ko}
고전적인 물리 기반 형상 제어는 외력으로 물질을 목표로 이끈다~\citep{treuille2003keyframe,mcnamara2004adjoint,fattal2004targetdriven,chen2024eigenfluids,li2025flowmapadjoint}.
우리의 선택을 시험하기 위해 $\delta\mathbf{F}_c$를 입자와 단계마다의 최적화된 외력으로 바꾸고 나머지는 모두 그대로 두었다.
그냥 힘은 매 창 운동량을 넣고, 목적 함수가 이를 다시 끌어와야 한다. 힘을 알짜 힘과 토크 0으로 투영하면 그 운동량은 정확히
없어지지만 형상의 부족은 없어지지 않는다. 두 힘 갈래 모두 응력 제어 기하 에너지의 몇 배와 더 나쁜 표면으로 끝나며, 더 정지한
채 끝나지도 않는다. 응력 제어는 투영 없이 정확한 운동량을 주고, 같은 목적 함수로 운동량 없는 힘이 도달하지 못하는 더 나은
끝을 준다.

\subsection{다른 방법과의 비교}
\label{ch4:sec_comparison_ko}
다섯 종류, 아홉 방법과 비교한다. 각 방법은 같은 구에서 같은 여섯 목표로 돌리고 같은 판독기로 읽는다.
\ref{CH3_SHAPEMORPHING}장의 F-제어~\citep{11088224}, Buonomo 등의 신경 암시적 모핑~\citep{buonomo2025vpnsm},
NISE~\citep{novello2023nise}, ISD~\citep{sang2025isd}, 4Deform~\citep{sang2025fourdeform}, 가능한 한 강체인
보간~\citep{alexa2000arap}, 최적 수송 변위 보간~\citep{mccann1997convexity,bonneel2011displacement}, 그리고 Large
Steps~\citep{nicolet2021largesteps}와 FlexiCubes~\citep{shen2023flexicubes}를 쓰는 렌더 전용 메시 최적화이다. 표~\ref{ch4:tab_methods_ko}와
부록~\ref{app:ch4_comparison_ko}에 전체 판독이 있다.

어느 메시에서도 다른 방법의 실루엣이 우리보다 위에 있지 않으며, 가장 뚜렷한 차이는 얇은 부분에 있다. 우리는 메시의 두 번째
샘플만큼만 얇은 목표 점을 덮지 못하고, 다른 모든 방법은 그 1.5배 이상을 남긴다. 그러나 방법들을 가장 크게 가르는 것은 가는
길이다. ISD는 끝을 맞추지만 도중에 몸체를 잃고, Large Steps는 genus 0 표면 하나를 유지하므로 손잡이를 결코 열 수 없으며,
FlexiCubes는 제어 없이 위상을 바꾸고, 수송만 쓰는 방법과 \ref{CH3_SHAPEMORPHING}장의 F-제어는 질량을 유지하지만 얇은 부분을
구름으로 찢는다. 우리는 정확한 질량과 거의 일정한 부피를 가진 한 몸체로 끝까지 남고, genus가 높은 것을 포함한 모든 목표에
도달한다. 각자의 끝 메시로 읽으면 가장 좋은 메시 방법들의 표면이 우리 입자 경로보다 날카롭지만, 목표의 완벽한 입자 샘플도 그
메시들에 비슷한 정도로 뒤지므로, 그 간극의 그 부분은 이 해상도에서 입자 몸체를 쓰는 비용이다. 렌더 전용 메시 방법들, ARAP,
F-제어, NISE는 우리보다 빠르고, ISD, 부피 보존 모핑, 4Deform은 느리다.

\subsection{실행 시간}
\label{ch4:sec_runtime_ko}
300k 모핑은 GPU 한 장에서 10분에서 30분쯤 걸린다(표~\ref{ch4:tab_runtime_ko}). 실행 길이는 정지 규칙이 정하고, 그 대부분은
도착 뒤의 느린 꼬리이다. 그 꼬리는 낭비가 아니다. 앞쪽 절반 동안 끝 척도가 계속 나아지고 몸체가 가라앉으며, 마지막 10분의
1만 잴 만한 것을 바꾸지 않는다.

\subsection{던지고 떨어뜨린 몸체}
\label{ch4:sec_throw_ko}
두 구를 서로를 향해 던지면, 반쯤 만들어진 채 공중에서 부딪히고, 떨어지고, 각각 bunny가 되어 착지하고 구른다
(그림~\ref{ch4:fig_throw_ko}). 장면은 부록~\ref{ch4:sec_scene_ko}의 여러 몸체 정의와 더 단단한 재질을 쓰며, 제어가 충격에
맞서 각 몸체의 형상을 붙잡도록 소성 동화를 끈다. 결과 실행에서 두 몸체는 모핑이 절반이 조금 안 되게 끝났을 때 만나고, 깨끗하게
떨어지며, 착지하여 바닥에서 튀어 오르고 구르며, 각자 자기 좌표계에서 목표에 가깝게 끝나고, 운동량은 모든 국면에서 유지된다.
몸체 하나가 떨어질 수도 있다. 바닥 위 높은 곳에서 놓은 구는 떨어지는 동안 bunny가 되어, 착지하고, 튀어 오른 뒤 머리 쪽으로
앞으로 넘어지는데, 몸체 좌표계 목적 함수가 이를 허용한다.

\begin{figure*}[tbp]
    \centering
    \includegraphics[width=\textwidth]{figs/ch4/fig_throw_scene.jpg}
    \caption{각각 bunny가 되는 던져진 두 몸체. \textmd{각 몸체가 자기 격자장과 자기 좌표계의 목적 함수를 가지므로, 깨끗하게
    떨어지고, 날고, 착지하고, 구른다.}}
    \label{ch4:fig_throw_ko}
\end{figure*}

% ============================================================================
\section{측정된 비용과 다음 단계}
\label{ch4:sec_costs_ko}
% ============================================================================

\textbf{정지도 지키는 렌더 가중치.} \ref{ch4:sec_write_ko}절은 렌더가 도착 뒤에도 권한을 유지하면 물리 목적 함수가 볼 수
없는 것, 곧 글자와 색 배치를 쓴다는 것을 보였다. 그러나 묶인 가중치는 몸체가 도착해도 사라지지 않으므로 실행이 멈출 때
몸체가 아직 움직인다. 열린 질문은 렌더가 정할 것이 있는 곳에서는 권한을 유지하고 몸체는 정지하게 하는 가중치의 정의이다.
시험 중인 정의는 렌더 항 자신의 측정된 바닥에 대한 초과분을 따라, 그림에만 있는 내용이 쓰이면 가중치가 사라지게 한다. 함께
시험하는 갈래는 물리 배치 항을 손실 셀 척도에서만 읽어 그보다 미세한 척도를 렌더에 맡긴다.
%% SWAP-IN: 바닥 기준 가중치와 "lambda 고정 + 손실 셀에서 읽는 배치 항" 갈래의 결과가 나오면 여기에 넣는다.

\textbf{일시적 부유물.} 초기 도착은 몸체에서 떨어진 입자 몇 개를 남기는데, 대부분 표시 표면 바깥으로 반 피치 서 있는 부조이다.
지속되는 것은 없고, 표시 표면에서는 따로 떨어진 조각을 이루지 않는다. 주된 원천은 $u$이다. 각 바깥층 단계를 이웃 가까이에
묶는 연결은 대부분을 없앴지만 질량 중심을 움직였고, 다음 형태는 최소 간격처럼 그 연결을 대칭으로 만든다.

\textbf{좁은 틈.} 3차 커널은 격자 속도를 약 두 셀에 걸쳐 평균하므로, 그보다 좁은 목표의 틈은 늦게까지 열리지 못하고 열릴 때
구슬로 찢어진다. $N$을 따르는 MPM 셀은 teapot 손잡이를 열고 경계를 날카롭게 하지만, dragon의 고리 대부분은 닫힌 채로 두고,
초기 표면을 구멍으로 채우며, dragon의 비용을 두 배 넘게 늘린다. 그래서 더 작은 셀은 파이프라인의 옵션이며, 더 많은 메시에서
같은 비교를 한 뒤 고른다.

\textbf{각운동량과 항력.} 주변 항력은 입자 운동량을 감쇠하지만 APIC의 아핀 부분은 감쇠하지 않는다. APIC 운동량 전체를 감쇠하는
항력은 시험 몸체에서 최대 각운동량을 몇 배 줄이고 형상과 정지를 유지하였다.

\textbf{길고 느린 꼬리.} 결과 프레임의 대부분은 도착 뒤에 온다. 몸체의 탄성 주기 동안 수송의 변위 보간을 따라 움직이는 목표는
변형을 시간에 걸쳐 퍼뜨리고, 운동 에너지의 최댓값과 초기 느슨한 입자를 줄이며, 렌더의 앞섬을 유지하고, 정지한 채 끝나지만,
몇 메시에서 그 끝이 고정 목표보다 조금 뒤진다. 다음 측정은 짧아진 재맞춤과 부족하게 공급된 끝부분을 구별하는 것이다.

\textbf{그림으로만 주어진 목표.} 3D 목표를 학습 실루엣의 visual hull로 두고 렌더 항만 진짜 메시를 보게 하면, 전체 방법은
홀드아웃 윤곽을 날카롭게 하지만 그림에만 있는 오목한 곳은 깎지 못한다. 정보는 그림에 있지만 표면 전체를 매끄럽게 바꿀 때만
손실이 내려간다. 다중 시점 shape-from-shading 하강은 오목한 곳을 제자리에서 찾으며, 다음 단계는 렌더 해상도에서 물체를
포함하는 hull에서 그곳을 깎는 것이다.

% ============================================================================
\section{장 요약}
\label{ch4:sec_summary_ko}
% ============================================================================

이 장은 \ref{CH3_SHAPEMORPHING}장의 미분 가능한 MLS-MPM 위에 세운 렌더 유도 모핑 파이프라인을 제시하였다. 몸체는 상한이 있는
응력 증분과 바깥층의 강체 성분 없는 오프셋으로 움직이는, 부피를 보존하는 고체이다. 몸체는 안정화 창에서 최적화되고, 각 창은
제어를 푼 뒤 수송 에너지, 안정 항과 정리 항, 그리고 자신의 외피 표면의 실루엣과 셰이딩으로 평가된다.

모핑은 정확한 질량과, 주변 항력을 빼면 운동량을 지키고, 정지한 채 끝나며, 구에는 없는 손잡이를 열고, 넓은 범위의 목표와 형상
쌍에 바꾸지 않고 적용되며, 다른 몸체와의 접촉과 중력에 참여한다. 다른 아홉 방법 중 어느 메시에서든 우리보다 실루엣이 높은
방법은 없고, 얇은 부분을 샘플링 바닥 수준으로 유지하는 것은 우리뿐이다. 렌더 항의 역할도 측정하였다. 기본 규칙에서 렌더는 몸체가 도착하는
동안 물질의 방향을 바꾸고, 그 뒤에는 가중치가 사라져 가는 물리 기울기를 따르므로 테두리를 다듬기만 하며, 끝을 일관되게 그러나
조금 정확한 목표에 더 가깝게 남긴다. 그리고 자신의 항에서 앞선 몫의 일부는 다른 형상이 아니라 더 매끈한 배열에서 온다.
도착 뒤에도 가중치를 유지하면 렌더는 목표 그림에만 있는 SIGGRAPH와 VANDERBILT라는 글자를 3D 목표가 매끈한 판에 새긴다.
같은 방식으로 그림이 빨갛게 칠한 bunny의 절반으로 빨간 물질을 옮기는데, 이 배치는 물리 목적 함수가 볼 수 없다. 둘 다 몸체의
정지를 대가로 치르며, 권한과 정지를 함께 지키는 렌더 가중치가 다음 단계이다.
